% Generated draft: requires mathematical and source review.
\documentclass[11pt]{article}
\usepackage[margin=25mm]{geometry}
\usepackage{fontspec}
\setmainfont{Latin Modern Roman}
\usepackage{amsmath,amssymb,mathtools,unicode-math}
\setmathfont{Latin Modern Math}
\usepackage{enumitem,needspace,xurl,hyperref}
\setlength{\parindent}{0pt}
\setlength{\parskip}{5pt}
\setlength{\emergencystretch}{3em}
\newcommand{\N}{\mathbb{N}}
\newcommand{\Z}{\mathbb{Z}}
\newcommand{\Q}{\mathbb{Q}}
\newcommand{\R}{\mathbb{R}}
\newcommand{\C}{\mathbb{C}}
\newcommand{\F}{\mathbb{F}}
\newcommand{\A}{\mathbb{A}}
\newcommand{\PP}{\mathbb{P}}
\newcommand{\E}{\mathbb{E}}
\newcommand{\eps}{\varepsilon}
\newcommand{\dd}{\,\mathrm{d}}
\begin{document}
\section*{UnsolvedMath ID: 11000108}
\textbf{Source problem code:} catalogue entry
\subsection{Definitions and mathematical statement}
Let \(G=\operatorname{Mod}(S_g)\) for a closed orientable surface in the source's intended genus range. The first question is whether there exists an endomorphism \(f:G\to G\) with
\[|f(G)|=\infty,\qquad [G:f(G)]=\infty.\]
The source conjectures that no such endomorphism exists. This concerns a homomorphism onto its own image, and neither surjectivity onto \(G\) nor injectivity of \(f\) is assumed.

Conditionally, if such \(f\) exists, the second question asks whether for some \(n\geq1\),
\[f|_{f^n(G)}:f^n(G)\to f^{n+1}(G)\text{ is injective}.\]
Equivalently,
\[\ker f\cap f^n(G)=\{1\},\quad\text{or}\quad\ker f^{n+1}=\ker f^n.\]
This is eventual injectivity on an iterated image, not a demand that the original map be injective. The genus range is not stated in this excerpt and matters for small-genus exceptional groups; no universal claim including every low genus is added.

The following context is about a different class of maps. Let \(\operatorname{Mod}_{g,p}^*\) be the extended mapping class group of a genus-\(g\) surface with \(p\) punctures, including orientation-reversing diffeomorphisms, and let \(\Gamma\leq\operatorname{Mod}_{g,p}^*\) have finite index. The source reports that an injective homomorphism \(\Gamma\to\operatorname{Mod}_{g,p}^*\) is the restriction of an automorphism of the ambient group, attributing this to Irmak [18,19] for \(g=2,p\geq2\) and for \(g\geq3\), and to Bell--Margalit [1] for \(g=0\). The genus-zero puncture/complexity restrictions and any exceptional cases are not reproduced, so this is a reported rigidity statement requiring exact hypotheses. It does not directly answer the noninjective endomorphism question above. The final sentence introduces another result that is absent.
\subsection{Short English statement}
Can an endomorphism of the closed-surface mapping class group have infinite image of infinite index? If so, must it become injective on some iterated image? The source conjectures nonexistence and separately cites finite-index embedding rigidity.
\subsection{Sources}
\begin{itemize}
\item Local UnsolvedMath registry, permanent ID \(11000108\), exact source code \texttt{catalogue entry}.
\item Farb, \emph{Problems on Mapping Class Groups and Related Topics}, Problem 2.2, PDF page 93: \url{https://www.math.uchicago.edu/~farb/papers/mcgbook.pdf}. Registry access 2026-07-29; pdf-font-structure extraction.
\item Irmak [18,19] and Bell--Margalit [1] are the supplied finite-index rigidity references; full entries and complete exceptional-case hypotheses are absent.
\end{itemize}
\end{document}
